
\documentclass[addpoints]{exam}
\usepackage{amsmath}
\usepackage{amssymb}
\usepackage{color}

% \printanswers

\newcommand{\event}{Characteristics of Functions}
\newcommand{\tournament}{}
\def\type{0}

\setlength\fillinlinelength{\textwidth}
\CorrectChoiceEmphasis{\color{red}\bfseries}
\SolutionEmphasis{\color{red}\bfseries}

\setlength\answerclearance{2pt}
\pagestyle{headandfoot}
\runningheadrule
\runningheader{\event}{\tournament}{\ifnum\type=1 Class Set Test \else Full Name:\hspace{1cm} \fi}
\runningfootrule
\runningfooter{}{Page \thepage\ of \numpages}{}

\begin{document}
\begin{center}
    {\huge Grade 11 Functions \\ \event
    \ifprintanswers \space Key
    \else \space \ifnum\type<2 Test \fi \ifnum\type=1 \textbf{(CLASS SET)} \fi
          \ifnum\type=2 Answer Sheet \fi
    \fi \\ \vspace{3mm}\tournament\vspace{1mm}}
\end{center}

\vspace{.25\textheight}

\begin{center}
    First Name: \ifprintanswers
    \underline{\hspace{3.5cm}}\textbf{KEY}\underline{\hspace{5.5cm}}\\\vspace{5mm}
    \else \underline{\hspace{10cm}}\\ \vspace{5mm} \fi
    Last Name: \underline{\hspace{10cm}}\\ \vspace{5mm}
\end{center}

\noindent\textbf{\underline{Directions:}}
\begin{itemize}
    \ifnum\type=1 \item \textbf{This test is a class set.} Do not write on this paper. \fi
    \item Show all work for full marks. Leave answers in exact form unless stated otherwise.
    \item The test is designed to be completed in 75 minutes.
\end{itemize}
\begin{center}
\textit{For grading use only}\\
\multirowgradetable{1}[pages]
\end{center}

\newpage
\begin{questions}

\section*{Part A --- Multiple Choice (15 marks)}
\vspace{5mm}

\question[1] Which of the following relations is \textbf{not} a function?
\begin{choices}
    \choice $\{(1,2),(2,3),(3,4)\}$
    \correctchoice $\{(1,2),(1,3),(2,4)\}$
    \choice $\{(0,0),(1,1),(2,4)\}$
    \choice $y = x^2$
\end{choices}

\question[1] If $f(x) = 3x^2 - 2x + 1$, then $f(-1) =$
\begin{choices}
    \choice $0$
    \choice $2$
    \correctchoice $6$
    \choice $-4$
\end{choices}

\question[1] The domain of $f(x) = \dfrac{1}{x - 3}$ is:
\begin{choices}
    \choice $\mathbb{R}$
    \choice $x > 3$
    \correctchoice $x \in \mathbb{R},\; x \neq 3$
    \choice $x \geq 3$
\end{choices}

\question[1] The range of $f(x) = x^2 + 2$ is:
\begin{choices}
    \choice $\mathbb{R}$
    \choice $y \geq 0$
    \correctchoice $y \geq 2$
    \choice $y > 2$
\end{choices}

\question[1] The graph of $y = f(x) + 3$ is obtained from $y = f(x)$ by:
\begin{choices}
    \choice Shifting 3 units right
    \choice Shifting 3 units left
    \correctchoice Shifting 3 units up
    \choice Shifting 3 units down
\end{choices}

\question[1] The graph of $y = f(x - 4)$ is obtained from $y = f(x)$ by:
\begin{choices}
    \choice Shifting 4 units up
    \correctchoice Shifting 4 units right
    \choice Shifting 4 units left
    \choice Reflecting in the $x$-axis
\end{choices}

\question[1] The graph of $y = -f(x)$ is a reflection of $y = f(x)$ in the:
\begin{choices}
    \choice $y$-axis
    \correctchoice $x$-axis
    \choice Line $y = x$
    \choice Origin
\end{choices}

\question[1] The graph of $y = f(-x)$ is a reflection of $y = f(x)$ in the:
\begin{choices}
    \correctchoice $y$-axis
    \choice $x$-axis
    \choice Line $y = x$
    \choice Origin
\end{choices}

\question[1] If $f(x) = 2x + 1$ and $g(x) = x^2$, then $(f \circ g)(3) =$
\begin{choices}
    \choice $7$
    \choice $25$
    \correctchoice $19$
    \choice $49$
\end{choices}

\question[1] The inverse of the function $f(x) = 2x - 6$ is:
\begin{choices}
    \choice $f^{-1}(x) = 2x + 6$
    \correctchoice $f^{-1}(x) = \dfrac{x+6}{2}$
    \choice $f^{-1}(x) = \dfrac{x}{2} - 6$
    \choice $f^{-1}(x) = \dfrac{1}{2x-6}$
\end{choices}

\question[1] For a function $f$, the graph of $f^{-1}$ is a reflection of $f$ over:
\begin{choices}
    \choice The $x$-axis
    \choice The $y$-axis
    \correctchoice The line $y = x$
    \choice The origin
\end{choices}

\question[1] The graph of $y = 3f(x)$ compared to $y = f(x)$ is:
\begin{choices}
    \choice Compressed vertically by factor 3
    \correctchoice Stretched vertically by factor 3
    \choice Stretched horizontally by factor 3
    \choice Shifted up 3 units
\end{choices}

\question[1] The graph of $y = f(2x)$ compared to $y = f(x)$ is:
\begin{choices}
    \choice Stretched horizontally by factor 2
    \correctchoice Compressed horizontally by factor $\tfrac{1}{2}$
    \choice Stretched vertically by factor 2
    \choice Shifted right 2 units
\end{choices}

\question[1] If $f(x) = x + 1$ and $g(x) = x^2 - 1$, then $(f \cdot g)(x) =$
\begin{choices}
    \choice $x^2 + x$
    \correctchoice $x^3 + x^2 - x - 1$
    \choice $x^2 - x - 2$
    \choice $x^3 - x$
\end{choices}

\question[1] The range of $f(x) = \sqrt{x - 1}$ is:
\begin{choices}
    \choice $x \geq 1$
    \choice All real numbers
    \correctchoice $y \geq 0$
    \choice $y > 0$
\end{choices}


\section*{Part B --- Short Answer (33 marks)}

\question Let $f(x) = x^2 - 4x + 3$ and $g(x) = 2x - 1$.
\begin{parts}
    \part[2] Find $(f + g)(x)$ and $(f - g)(x)$.
    \part[2] Find $(f \circ g)(x)$. Simplify fully.
    \part[2] Find $(g \circ f)(2)$.
    \part[2] Find all values of $x$ where $f(x) = g(x)$.
\end{parts}
\begin{solution}[10cm]
\textbf{(a)} $(f+g)(x) = x^2 - 4x + 3 + 2x - 1 = \mathbf{x^2 - 2x + 2}$

$(f-g)(x) = x^2 - 4x + 3 - (2x-1) = \mathbf{x^2 - 6x + 4}$

\textbf{(b)} $(f\circ g)(x) = f(2x-1) = (2x-1)^2 - 4(2x-1) + 3$
$= 4x^2 - 4x + 1 - 8x + 4 + 3 = \mathbf{4x^2 - 12x + 8}$

\textbf{(c)} $f(2) = 4 - 8 + 3 = -1$;\quad $(g\circ f)(2) = g(-1) = 2(-1)-1 = \mathbf{-3}$

\textbf{(d)} $x^2 - 4x + 3 = 2x - 1 \Rightarrow x^2 - 6x + 4 = 0$
\[
x = \frac{6 \pm \sqrt{36-16}}{2} = \frac{6 \pm \sqrt{20}}{2} = 3 \pm \sqrt{5}
\]
$\mathbf{x = 3 + \sqrt{5}}$ or $\mathbf{x = 3 - \sqrt{5}}$
\end{solution}

\question For the function $f(x) = \sqrt{x + 4}$:
\begin{parts}
    \part[1] State the domain and range of $f$.
    \part[3] Find $f^{-1}(x)$. State its domain and range.
    \part[2] Verify algebraically that $f(f^{-1}(x)) = x$ for all $x$ in the domain
    of $f^{-1}$.
    \part[2] If $g(x) = x - 1$, find the domain of $\dfrac{f}{g}(x) = \dfrac{\sqrt{x+4}}{x-1}$.
\end{parts}
\begin{solution}[9cm]
\textbf{(a)} Domain of $f$: $x \geq -4$;\quad Range of $f$: $y \geq 0$

\textbf{(b)} Let $y = \sqrt{x+4}$. Swap: $x = \sqrt{y+4} \Rightarrow x^2 = y+4 \Rightarrow y = x^2 - 4$.
\[
f^{-1}(x) = x^2 - 4
\]
Domain of $f^{-1}$: $x \geq 0$ (range of $f$);\quad Range of $f^{-1}$: $y \geq -4$

\textbf{(c)} $f(f^{-1}(x)) = f(x^2-4) = \sqrt{(x^2-4)+4} = \sqrt{x^2} = |x| = x$
for $x \geq 0$ \checkmark

\textbf{(d)} Need $x + 4 \geq 0$ (so $x \geq -4$) and $x - 1 \neq 0$ (so $x \neq 1$).
Domain: $\mathbf{x \geq -4,\; x \neq 1}$
\end{solution}

\question Describe the transformations applied to $y = f(x)$ to obtain each function,
then state the effect on domain and range if the original domain is $[0, 4]$ and
range is $[1, 5]$.
\begin{parts}
    \part[3] $y = 2f(x - 3) + 1$
    \part[3] $y = -f(2x)$
\end{parts}
\begin{solution}[9cm]
\textbf{(a)} Transformations (in order): horizontal shift right 3, vertical stretch by
factor 2, vertical shift up 1.

New domain: $[0+3,\; 4+3] = [3,\; 7]$

New range: $[2(1)+1,\; 2(5)+1] = [3,\; 11]$

\textbf{(b)} Transformations: horizontal compression by factor $\tfrac{1}{2}$
(i.e.\ $x$-values halved), reflection in $x$-axis.

New domain: $[0/2,\; 4/2] = [0,\; 2]$

New range: $[-5,\; -1]$ (range negated by reflection)
\end{solution}

\question A function is defined as $f(x) = \dfrac{1}{2}(x + 1)^2 - 3$.
\begin{parts}
    \part[2] Describe the transformations from $y = x^2$ to $y = f(x)$.
    \part[2] State the vertex, domain, and range of $f$.
    \part[3] Find $f^{-1}(x)$. Is $f^{-1}$ a function? If not, how can you restrict
    the domain of $f$ so that $f^{-1}$ is a function?
\end{parts}
\begin{solution}[9cm]
\textbf{(a)} From $y = x^2$: shift left 1 unit, vertical compression by $\tfrac{1}{2}$,
shift down 3 units.

\textbf{(b)} Vertex: $(-1, -3)$;\quad Domain: $\mathbb{R}$;\quad Range: $y \geq -3$

\textbf{(c)} $y = \tfrac{1}{2}(x+1)^2 - 3$. Swap $x$ and $y$:
$x = \tfrac{1}{2}(y+1)^2 - 3 \Rightarrow x + 3 = \tfrac{1}{2}(y+1)^2$
$\Rightarrow (y+1)^2 = 2(x+3) \Rightarrow y = -1 \pm\sqrt{2(x+3)}$

$f^{-1}(x) = -1 \pm\sqrt{2(x+3)}$, defined for $x \geq -3$.

\textbf{Not a function} (fails vertical line test — $\pm$ gives two values).
Restrict domain of $f$ to $x \geq -1$ (right branch) to get
$f^{-1}(x) = -1 + \sqrt{2(x+3)}$, or to $x \leq -1$ (left branch) for the
$-$ version.
\end{solution}

\end{questions}
\end{document}
